% Importado desde: 2026-03-28_Curvatura_Berry_y_Carga_Quiral_Weyl.tex
\chapter{Derivacion Pedagogica: --- curvatura de Berry y carga quiral en un par Weyl}
\label{ch:berry-carga-quiral-weyl}






% verified-by: validators/chirality.py::test_canonical_weyl_det
% verified-by: validators/chirality.py::test_opposite_weyl_opposite_chirality
% verified-by: validators/chirality.py::test_parity_flips_chirality
% verified-by: validators/chirality.py::test_three_sigma_trace_levi_civita

\section{Objetivo}
En el capitulo anterior obtuvimos, de manera controlada, un par de nodos Weyl separados en el espacio recíproco. Ahora queremos entender su lectura geométrica.

La pregunta central de este paso es:
\begin{displayquote}
\textbf{que significa geometricamente decir que un nodo Weyl tiene quiralidad \(+1\) o \(-1\), y como se expresa eso en terminos de curvatura de Berry y flujo topologico?}
\end{displayquote}

Este paso es importante porque traduce el lenguaje algebraico de quiralidad a un lenguaje geométrico y topológico. Eso nos permite entender:
\begin{enumerate}[label=\arabic*.]
    \item por que un nodo Weyl actua como una fuente o sumidero;
    \item por que la carga quiral se conserva globalmente;
    \item por que un par Weyl aislado infrarrojo sigue obedeciendo la contabilidad global de la zona de Brillouin.
\end{enumerate}

\section{El modelo local de un nodo Weyl}
Trabajamos primero con un solo nodo Weyl idealizado, centrado en \(\bm{q}=0\):
\begin{equation}
H_\chi(\bm{q})=
\chi\, v\, \bm{\sigma}\cdot \bm{q},
\qquad
\chi=\pm 1.
\label{c10:eq:weyl}
\end{equation}

Aqui:
\begin{enumerate}[label=\arabic*.]
    \item \(v\) es la velocidad efectiva;
    \item \(\bm{q}\) mide la distancia al nodo;
    \item \(\chi\) es la quiralidad.
\end{enumerate}

El espectro es
\begin{equation}
E_\pm(\bm{q})=\pm v\abs{\bm{q}}.
\end{equation}

Las autofunciones dependen de la direccion de \(\bm{q}\), no solo de su modulo. Esa dependencia angular es justamente la puerta de entrada a la conexion de Berry.

\section{Por que aparece la geometria de Berry}
Cuando un Hamiltoniano depende suavemente de parametros, sus autovectores pueden arrastrar una geometria intrinseca en el espacio de esos parametros.

Aqui el espacio de parametros es el espacio recíproco local, coordinado por \(\bm{q}\). Por tanto:
\begin{enumerate}[label=\arabic*.]
    \item el estado de banda \(\ket{u_n(\bm{q})}\) cambia al movernos en \(\bm{q}\);
    \item ese cambio define una conexion;
    \item la rotacional de esa conexion define una curvatura.
\end{enumerate}

La conexion de Berry de una banda \(n\) es
\begin{equation}
\bm{\mathcal{A}}_n(\bm{q})
=
i\bra{u_n(\bm{q})}\bm{\nabla}_{\bm{q}}\ket{u_n(\bm{q})},
\label{c10:eq:berryconn}
\end{equation}
y la curvatura de Berry es
\begin{equation}
\bm{\Omega}_n(\bm{q})
=
\bm{\nabla}_{\bm{q}}\times \bm{\mathcal{A}}_n(\bm{q}).
\label{c10:eq:berrycurv}
\end{equation}

\subsection{Interpretacion minima}
La conexion depende de la fase local elegida para los autovectores, pero la curvatura es invariante de gauge. Por eso la cantidad fisica natural es \(\bm{\Omega}_n\), no \(\bm{\mathcal{A}}_n\) por separado.

\begin{figure}[h!]
\centering
\begin{tikzpicture}[
    box/.style={draw, rounded corners, thick, align=center, minimum width=3.8cm, minimum height=1.3cm},
    arr/.style={-{Latex[length=3mm]}, thick}
]
\node[box, fill=blue!7] (a) at (0,0) {autovectores\\\(\ket{u_n(\bm{q})}\)};
\node[box, fill=yellow!10] (b) at (5,0) {conexion\\\(\bm{\mathcal{A}}_n\)};
\node[box, fill=green!7] (c) at (10,0) {curvatura\\\(\bm{\Omega}_n\)};
\node[box, fill=red!7] (d) at (15,0) {flujo y\\carga quiral};
\draw[arr] (a) -- (b);
\draw[arr] (b) -- (c);
\draw[arr] (c) -- (d);
\end{tikzpicture}
\caption{La secuencia conceptual del capitulo.}
\end{figure}

\section{Autovectores del Hamiltoniano Weyl}
Es util escribir
\begin{equation}
\bm{q} = q(\sin\theta\cos\phi,\sin\theta\sin\phi,\cos\theta).
\end{equation}

Para \(\chi=+1\), una eleccion conveniente de autovectores normalizados es
\begin{equation}
\ket{u_+(\theta,\phi)}=
\begin{pmatrix}
\cos(\theta/2)\\[0.4em]
e^{i\phi}\sin(\theta/2)
\end{pmatrix},
\qquad
\ket{u_-(\theta,\phi)}=
\begin{pmatrix}
-e^{-i\phi}\sin(\theta/2)\\[0.4em]
\cos(\theta/2)
\end{pmatrix}.
\label{c10:eq:eigenstates}
\end{equation}

Estas funciones ya muestran algo importante: el estado depende solo de la direccion de \(\bm{q}\), es decir, del punto sobre la esfera en espacio de momentos.

\section{Calculo directo de la conexion de Berry}
Tomemos la banda positiva \(\ket{u_+}\). En coordenadas esfericas, la unica componente no trivial de la conexion sera la azimutal.

\subsection{Componente polar}
Calculamos
\begin{equation}
\mathcal{A}_\theta
=
i\bra{u_+}\partial_\theta u_+\rangle.
\end{equation}
Usando \eqref{c10:eq:eigenstates} se obtiene
\begin{equation}
\mathcal{A}_\theta=0.
\end{equation}

\subsection{Componente azimutal}
Ahora
\begin{equation}
\partial_\phi \ket{u_+}
=
\begin{pmatrix}
0\\
i e^{i\phi}\sin(\theta/2)
\end{pmatrix}.
\end{equation}
Entonces
\begin{align}
\mathcal{A}_\phi
&=
i\bra{u_+}\partial_\phi u_+\rangle
\\
&=
i\left[
\cos(\theta/2)\cdot 0
+
e^{-i\phi}\sin(\theta/2)\cdot i e^{i\phi}\sin(\theta/2)
\right]
\\
&=
i\left[i\sin^2(\theta/2)\right]
\\
&=
-\sin^2(\theta/2).
\end{align}

Usando
\begin{equation}
\sin^2(\theta/2)=\frac{1-\cos\theta}{2},
\end{equation}
podemos escribir
\begin{equation}
\mathcal{A}_\phi = -\frac{1-\cos\theta}{2}.
\label{c10:eq:Aphi}
\end{equation}

Esta expresion tiene la forma típica del potencial de un monopolo magnético en una carta local.

\section{Curvatura de Berry}
La curvatura en coordenadas esfericas sale de
\begin{equation}
\Omega_r
=
\frac{1}{q^2\sin\theta}
\partial_\theta \mathcal{A}_\phi.
\end{equation}
Usando \eqref{c10:eq:Aphi},
\begin{equation}
\partial_\theta \mathcal{A}_\phi
=
-\frac{1}{2}\sin\theta.
\end{equation}
Entonces
\begin{equation}
\Omega_r
=
-\frac{1}{2q^2}.
\end{equation}

Es decir, para la banda positiva del nodo con \(\chi=+1\),
\begin{equation}
\bm{\Omega}_+(\bm{q})
=
-\frac{1}{2}\frac{\hat{\bm{q}}}{q^2}
=
-\frac{\bm{q}}{2\abs{\bm{q}}^3}.
\label{c10:eq:omega-plus}
\end{equation}

Para la banda negativa el signo cambia:
\begin{equation}
\bm{\Omega}_-(\bm{q})
=
\frac{\bm{q}}{2\abs{\bm{q}}^3}.
\label{c10:eq:omega-minus}
\end{equation}

\subsection{Con quiralidad general}
Si el nodo tiene quiralidad \(\chi=\pm 1\), la expresion se escribe de forma compacta como
\begin{equation}
\boxed{
\bm{\Omega}_{\text{val}}(\bm{q})
=
\chi\,\frac{\bm{q}}{2\abs{\bm{q}}^3},
\qquad
\bm{\Omega}_{\text{cond}}(\bm{q})
=
-\chi\,\frac{\bm{q}}{2\abs{\bm{q}}^3}.
}
\label{c10:eq:omegageneral}
\end{equation}

Aqui `val` y `cond` denotan banda de valencia y de conduccion.

\section{Lectura geométrica: un monopolo en espacio recíproco}
La forma \eqref{c10:eq:omegageneral} es exactamente la de un campo radial de monopolo:
\begin{equation}
\bm{\Omega}(\bm{q}) \propto \frac{\bm{q}}{\abs{\bm{q}}^3}.
\end{equation}

Por eso se dice que un nodo Weyl es una fuente o sumidero topológico:
\begin{enumerate}[label=\arabic*.]
    \item \(\chi=+1\): fuente de flujo;
    \item \(\chi=-1\): sumidero de flujo.
\end{enumerate}

\begin{figure}[h!]
\centering
\begin{tikzpicture}[scale=1.2]
\draw[thick] (0,0) circle (2);
\fill[blue!70!black] (0,0) circle (1.5pt);
\foreach \ang in {0,30,...,330} {
    \draw[-{Latex[length=2.5mm]}, blue!70!black, thick] (0,0) -- ({1.8*cos(\ang)},{1.8*sin(\ang)});
}
\node[blue!70!black] at (0,-2.45) {\(\chi=+1\): fuente};

\begin{scope}[xshift=8cm]
\draw[thick] (0,0) circle (2);
\fill[red!70!black] (0,0) circle (1.5pt);
\foreach \ang in {0,30,...,330} {
    \draw[-{Latex[length=2.5mm]}, red!70!black, thick] ({1.8*cos(\ang)},{1.8*sin(\ang)}) -- (0,0);
}
\node[red!70!black] at (0,-2.45) {\(\chi=-1\): sumidero};
\end{scope}
\end{tikzpicture}
\caption{Lectura geométrica de la quiralidad como flujo radial de curvatura de Berry.}
\end{figure}

\section{Flujo y carga topologica}
La cantidad topológica relevante es el flujo de la curvatura de Berry a través de una esfera \(S^2\) que encierra el nodo:
\begin{equation}
\Phi = \oint_{S^2} \bm{\Omega}\cdot d\bm{S}.
\label{c10:eq:flux}
\end{equation}

Tomando la banda de valencia y usando \eqref{c10:eq:omegageneral},
\begin{equation}
\bm{\Omega}=
\chi\,\frac{\hat{\bm{q}}}{2q^2},
\qquad
d\bm{S}= \hat{\bm{q}}\, q^2\sin\theta\, d\theta\, d\phi.
\end{equation}
Entonces
\begin{align}
\Phi
&=
\int_0^{2\pi}\int_0^\pi
\chi\,\frac{1}{2q^2}
q^2\sin\theta\, d\theta\, d\phi
\\
&=
\frac{\chi}{2}
\int_0^{2\pi}d\phi
\int_0^\pi \sin\theta\, d\theta
\\
&=
\frac{\chi}{2}(2\pi)(2)
\\
&=
2\pi \chi.
\label{c10:eq:fluxresult}
\end{align}

El numero de Chern asociado es
\begin{equation}
C = \frac{1}{2\pi}\Phi = \chi.
\label{c10:eq:chern}
\end{equation}

Esta ecuacion es el corazon topologico del capitulo:
\begin{displayquote}
\textbf{la quiralidad de un nodo Weyl es exactamente su carga monopolar de curvatura de Berry, o equivalentemente el numero de Chern del cascaron esferico que lo rodea.}
\end{displayquote}

\section{Volver al par Weyl que construimos}
En el capitulo anterior obtuvimos dos nodos en
\begin{equation}
\bm{q}_\pm = (0,0,\pm q_0),
\qquad
q_0=\frac{\sqrt{b^2-m^2}}{v},
\end{equation}
con quiralidades opuestas:
\begin{equation}
\chi_+=+1,
\qquad
\chi_-=-1.
\end{equation}

Por tanto:
\begin{enumerate}[label=\arabic*.]
    \item el nodo en \(+q_0\) emite flujo \(+2\pi\);
    \item el nodo en \(-q_0\) absorbe flujo \(-2\pi\);
    \item la suma total del flujo es cero.
\end{enumerate}

\subsection{Compatibilidad con Nielsen--Ninomiya}
Esto encaja perfectamente con lo ya discutido. Aunque en el infrarrojo vemos un par Weyl limpio, la carga total sigue cancelandose:
\begin{equation}
\chi_+ + \chi_- = 0.
\end{equation}

No hemos violado Nielsen--Ninomiya. Lo que hicimos fue aislar un par infrarrojo bien controlado, no crear una carga quiral neta global en la red.

\section{Lectura por cortes bidimensionales}
Hay otra manera muy util de ver la misma fisica. Fijemos \(q_z\) y miremos el sistema como una familia de problemas \(2D\). Entonces:
\begin{enumerate}[label=\arabic*.]
    \item para valores de \(q_z\) entre ambos nodos, el corte \(2D\) tiene un numero de Chern no trivial;
    \item fuera de esa region, el numero de Chern cambia;
    \item el salto ocurre precisamente al cruzar un nodo Weyl.
\end{enumerate}

En otras palabras, cada nodo actua como punto donde el invariante topologico de los cortes \(2D\) cambia en una unidad.

\begin{figure}[h!]
\centering
\begin{tikzpicture}[
    box/.style={draw, rounded corners, thick, align=center, minimum width=3.8cm, minimum height=1.3cm},
    arr/.style={-{Latex[length=3mm]}, thick}
]
\node[box, fill=blue!7] (a) at (0,0) {cortes \(2D\)\\para \(q_z< -q_0\)};
\node[box, fill=yellow!10] (b) at (5,0) {cruce del nodo\\\(-q_0\)};
\node[box, fill=green!7] (c) at (10,0) {region entre nodos\\Chern distinto};
\node[box, fill=yellow!10] (d) at (15,0) {cruce del nodo\\\(+q_0\)};
\draw[arr] (a) -- (b);
\draw[arr] (b) -- (c);
\draw[arr] (c) -- (d);
\end{tikzpicture}
\caption{Los nodos Weyl marcan saltos topológicos en los cortes bidimensionales.}
\end{figure}

\section{Que hemos aprendido de verdad}
Conviene resumir lo ganado con cuidado.

\subsection*{Primero}
La quiralidad no es solo el signo del determinante de la matriz de velocidades. Tambien es una carga topologica medible por flujo de Berry.

\subsection*{Segundo}
La curvatura de Berry de una banda cerca de un nodo Weyl tiene forma monopolar:
\begin{equation}
\bm{\Omega}(\bm{q})\sim \chi \frac{\bm{q}}{2\abs{\bm{q}}^3}.
\end{equation}

\subsection*{Tercero}
El flujo a traves de una esfera alrededor del nodo vale \(2\pi\chi\), y el numero de Chern correspondiente es \(\chi\).

\subsection*{Cuarto}
En el par Weyl que construimos antes, las dos cargas son opuestas y la suma global se cancela, como debe ocurrir en una red completa.

\section{Por que este paso importa para el proyecto}
Ahora la escalera conceptual queda todavia mas firme:
\begin{enumerate}[label=\arabic*.]
    \item construimos una dinamica discreta;
    \item controlamos los dobladores;
    \item aislamos un sector Dirac ligero;
    \item lo desdoblamos en un par Weyl controlado;
    \item entendimos la quiralidad de ese par como una carga topologica.
\end{enumerate}

Eso es importante porque el programa ya no solo reproduce una ecuacion efectiva: tambien reproduce la geometria topologica que distingue a los semimetales de Weyl.

\section{Mapa del siguiente paso}
Ahora los dos caminos mas naturales son:
\begin{enumerate}[label=\arabic*.]
    \item estudiar como esta estructura topologica se relaciona con estados de borde y arcos de Fermi;
    \item volver a la pregunta de simetria efectiva y analizar como conviven los generadores tipo Lorentz con una geometria de Berry no trivial.
\end{enumerate}

\section{Resumen final}
Partiendo del Hamiltoniano local de Weyl, calculamos paso a paso la conexion de Berry, obtuvimos la curvatura correspondiente y vimos que toma exactamente la forma de un campo monopolar en espacio recíproco. El flujo a traves de una esfera que rodea al nodo vale
\[
\Phi = 2\pi \chi,
\]
de modo que la quiralidad \(\chi\) coincide con la carga topologica del nodo.

Aplicado al par Weyl construido en el capitulo anterior, esto significa que uno de los nodos es fuente y el otro sumidero de curvatura de Berry. La suma total sigue siendo cero, en acuerdo con Nielsen--Ninomiya. Con esto ya no solo tenemos un par Weyl controlado algebraicamente: tambien entendemos su significado geométrico y topológico.
